\section{Status and Plan}
\label{sec:status}

\subsection{Work completed}
The pipeline of \cref{sec:method} has been built end to end, on TIMME, on the Pokec
district slice and on the synthetic generator. Each stage is a separate Python module
with its own self-check on a small planted case with a known answer:
\begin{itemize}
  \item \textbf{Topology view}: repeated Leiden runs, the reference partition $t$ and
        the set-aside small communities (\cref{sec:method-views}).
  \item \textbf{Content view and scores}: structural soft membership by label
        spreading, the cross-fitted text classifier, and the misfit, disagreement and
        baseline scores (\cref{sec:method-scores}). Diagnostics that check whether a
        partition or an embedding just sorts nodes by degree were added at the same
        time.
  \item \textbf{Detection checks}: the synthetic benchmark with planted conflicts and
        the TIMME party-misfit check with its shuffled-text null
        (\cref{sec:method-detect}).
  \item \textbf{Node blocking}: content-blind and content-aware Independent Cascade,
        all node-blocking strategies including the gateway score
        (\cref{sec:method-block}).
  \item \textbf{Regime map}: the node-level sweep over the synthetic grid
        (\cref{tab:synthetic}), and the placement of TIMME on that map using its
        estimated agreement effect and the degree position of its mismatched users.
  \item \textbf{Link-level intervention}: link friction with the structural and
        content-aware link scores.
\end{itemize}
The earlier GNN stage (\cref{sec:method-evolution}) was built and trained, and then
replaced; it is not part of the current pipeline.

\subsection{Work remaining}
\begin{itemize}
  \item \textbf{Edge-level regime map.} The regime map currently covers node blocking
        only. The same sweep is to be run for link-level intervention, followed by a
        write-up of the overall argument in the light of both maps.
  \item \textbf{Full results for the final report.} All detection, blocking and regime
        outputs will be reported, with their baselines and uncertainty, in the final
        report.
  \item \textbf{Full Pokec.} The project plan develops on the Pokec district slice and
        keeps the full Pokec graph for the final run. The final run will use the
        largest strongly connected component of the directed Pokec friendship graph.
  \item \textbf{Supervisor check-in.} The narrowed claim (\cref{sec:related-gap}) is to
        be discussed with the supervisor.
  \item \textbf{Active-learning extension} (\cref{sec:method-al}). Whether this is
        done depends on the supervisor's decision.
\end{itemize}

\subsection{Schedule}
\Cref{tab:schedule} lists the remaining steps in order. All work is planned to finish by 30 November 2026.

\begin{table}[ht]
\centering
\caption{Plan up to the final report.}
\label{tab:schedule}
\small
\begin{tabularx}{\linewidth}{@{}l X l@{}}
\toprule
Step & Work & Target date \\
\midrule
1 & Supervisor check-in on the narrowed claim and the active-learning extension & 21 Oct 2026 \\
2 & Edge-level regime map & 31 Oct 2026 \\
3 & Write-up of the overall argument from the node- and link-level maps & 7 Nov 2026 \\
4 & Full-Pokec run (largest strongly connected component) & 14 Nov 2026 \\
5 & Active-learning extension, if approved & 23 Nov 2026 \\
6 & Final report & 30 Nov 2026 \\
\bottomrule
\end{tabularx}
\end{table}
